\exercisesection{\S~5}

\begin{enumerate}
\def\labelenumi{\arabic{enumi})}
\item
  \(\operatorname{Aut}_{0}(\mathfrak{g})\) 在 \(\operatorname{Aut}(\mathfrak{g})\) 中的指数是有限的。
\item
  有 \(\mathrm{Aut}_e(\mathfrak{g}) = \mathrm{Aut}(\mathfrak{g})\)，当 \(\mathfrak{g}\) 可分裂、单且属于 \(\mathrm{G}_2, \mathrm{F}_4\) 或 \(\mathrm{E}_8\) 型时。
\item
  令 \(\mathfrak{h}\) 为 \(\mathfrak{g}\) 的分裂嘉当子代数，\(\mathfrak{b}\) 为 \((\mathfrak{g},\mathfrak{h})\) 的 Borel 子代数，\(\mathfrak{n} = [\mathfrak{b},\mathfrak{b}]\)，且 \(N = \exp \operatorname{ad}_{\mathfrak{g}}\mathfrak{n}\)。则
\end{enumerate}

\[
\operatorname{Aut} (\mathfrak {g}) = \mathrm{N}. \operatorname{Aut} (\mathfrak {g}, \mathfrak {h}). \mathrm{N}.
\]

(设 \(s \in \operatorname{Aut}(\mathfrak{g})\)。将§3第3节的命题10应用于 \(\mathfrak{b} \cap s(\mathfrak{b})\)，然后应用第VII章§3第4节的定理3。）

\begin{enumerate}
\def\labelenumi{\arabic{enumi})}
\setcounter{enumi}{3}
\tightlist
\item
  令 \(\mathfrak{h}\) 为 \(\mathfrak{g}\) 的嘉当子代数，令 \(s\) 为 \(\operatorname{Aut}(\mathfrak{g},\mathfrak{h})\) 中的一个元素，满足 \(sH \neq H\) 对所有非零 \(H\) 属于 \(\mathfrak{h}\) 成立。证明 \(s\) 的阶有限。（化归到 \(\mathfrak{h}\) 分裂的情形，并选取整数 \(n \geq 1\) 使得 \(\varepsilon(s)^n = 1\)。于是存在 \(\varphi \in \mathrm{T}_{\mathrm{Q}}\) 使得 \(f(\varphi) = s^n\)。令 \(\sigma\) 为 \(s|\mathfrak{h}\) 的转置。证明 \(1 + \sigma + \sigma^2 + \cdots + \sigma^{n-1} = 0\)，并由此推出 \(s^{n^2} = 1\)。）
\end{enumerate}

¶ 5) a) 令 \(a \in \operatorname{Aut}(\mathfrak{g})\)，令 \(\mathfrak{n}\) 为 \(a - 1\) 的幂零空间。证明下列条件等价：

\begin{enumerate}
\def\labelenumi{(\roman{enumi})}
\item
  \(\operatorname{Ker}(a - 1)\) 是 \(\mathfrak{g}\) 的嘉当子代数。
\item
  \(\mathfrak{n}\) 是 \(\mathfrak{g}\) 的嘉当子代数，且 \(a\in \mathrm{Aut}_0(\mathfrak{g})\)。
\item
  \(\dim\mathfrak{n}=\operatorname{rk}(\mathfrak{g})\)，且 \(a\in\operatorname{Aut}_{0}(\mathfrak{g})\)。
\end{enumerate}

\begin{enumerate}
\def\labelenumi{\alph{enumi})}
\setcounter{enumi}{1}
\item
  从现在起假设 k 是代数闭的。设 V 是向量空间，R 是 V 中的根系，\(T_{Q}\) 是群 \(\operatorname{Hom}(Q(R), k^{*})\)，n 是一个 \(\geq 1\) 的整数，而 \(T_{n}\) 是 \(T_{Q}\) 的子群，由阶整除 n 的元素组成。~设 \(\zeta\) 是本原 \(n\) 次单位根，属于 \(k\)。对于所有 \(H \in \mathrm{P}(\mathrm{R}^{\vee})\)，令 \(\psi(H)\) 为 \(\gamma \mapsto \zeta^{\gamma(H)}\) 所表示的 \(\mathrm{T}_{\mathrm{Q}}\) 中元素。映射 \(\psi\) 是从 \(\mathrm{P}(\mathrm{R}^{\vee})\) 到 \(\mathrm{T}_n\) 的同态，其核为 \(n\mathrm{P}(\mathrm{R}^{\vee})\)。令 \(t \in \mathrm{T}_n\)，令 \(\mathrm{C}\) 为 \(\mathrm{P}(\mathrm{R}^{\vee}) \otimes \mathbf{R}\) 中的一个单房。存在 \(w \in \mathrm{W}(\mathrm{R})\) 和 \(H \in \mathrm{P}(\mathrm{R}^{\vee})\)，使得 \(\frac{1}{n} H \in \overline{\mathrm{C}}\) 且 \(\psi(wH) = t\)。（使用第VI章§2第1节。）
\item
  令 \(\mathfrak{h}\) 为 \(\mathfrak{g}\) 的嘉当子代数，令 \(\mathbb{R}\) 为 \((\mathfrak{g},\mathfrak{h})\) 的根系。令 \(n,\zeta ,\psi\) 如 \(b\) ) 中，令 \(f\) 如第2节中。令 \(H\in \mathrm{P}(\mathrm{R}^{\vee})\)。在 \(\mathfrak{g}\) 中，\(f(\psi (H))\) 不变的元素集合为 \(\mathfrak{h}\oplus \sum_{\alpha \in \mathbb{R}'}\mathfrak{g}^\alpha\)，其中 \(\mathbb{R}'\) 是由 \(\alpha \in \mathbb{R}\) 且满足 \(\alpha (H)\in n\mathbf{Z}\) 的元素组成的集合；并且 \(f(\psi (H))\) 满足 \(a\) ) 中的条件，当且仅当 \(\frac{1}{n} H\) 属于一个单房。
\end{enumerate}

\(d\) ) 从现在起假设 \(\mathfrak{g}\) 是单的。令 \(\mathfrak{h}\) 和 \(\mathbb{R}\) 如 \(c\) ) 中，令 \((\alpha_{1},\ldots ,\alpha_{l})\) 为 \(\mathbb{R}\) 的一个基，令 \(h\) 为 \(\mathbb{R}\) 的 Coxeter 数，令 \(\zeta\) 为本原 \(h\) 次单位根，属于 \(k\)，并令 \(H\) 为 \(\mathfrak{h}\) 中满足 \(\alpha_{i}(H) = 1\)（\(i = 1,\dots,l\)）的元素。证明下列性质：

\begin{enumerate}
\def\labelenumi{(\roman{enumi})}
\item
  同态 \(\gamma \mapsto \zeta^{\gamma (H)}\) 从 Q(R) 到 \(k^{*}\) 定义了 \(\operatorname {Aut}(\mathfrak{g},\mathfrak{h})\) 中的一个元素，该元素满足 \(a\) ) 中的条件，且阶为 \(h\)。（使用 \(c\) )、第VI章§2的命题5和第VI章§1的命题31。）
\item
  \(\mathfrak{g}\) 的每个有限阶自同构都满足 \(a\) ) 中的条件，且其阶 \(\geq h\)。
\item
  \(\mathfrak{g}\) 的阶为 \(h\) 且满足 \(a\) ) 中条件的自同构在 \(\mathrm{Aut}_e(\mathfrak{g})\) 中构成一个共轭类。（使用第3节的命题5。）
\end{enumerate}

\begin{enumerate}
\def\labelenumi{\alph{enumi})}
\setcounter{enumi}{4}
\tightlist
\item
  令 \(\mathfrak{h}\) 和 \(\mathrm{R}\) 如 \(c\) ) 中，令 \(w\) 为 \(\mathrm{W(R)}\) 中的 Coxeter 变换。令 \(s \in \operatorname{Aut}(\mathfrak{g}, \mathfrak{h})\) 满足 \(\varepsilon(s) = w\)。证明 \(s\) 满足 \(a\) ) 中的条件，且阶为 \(h\)。（使用第VI章§1的命题33、第V章§6第2节和第VII章§4的命题9。）
\end{enumerate}

\(f)\) 若 \(s \in \operatorname{Aut}(\mathfrak{g})\)，则下列条件等价：

\begin{enumerate}
\def\labelenumi{(\roman{enumi})}
\item
  \(s\) 满足 \(a\) ) 中的条件，且阶为 \(h\)；
\item
  存在一个 \(\mathfrak{h}\)，它是 \(\mathfrak{g}\) 的在 \(s\) 下稳定的嘉当子代数，使得 \(s|\mathfrak{h}\) 是 \((\mathfrak{g},\mathfrak{h})\) 的 Weyl 群中的 Coxeter 变换。（使用 \(d\) ) 和 \(e\) )。）
\end{enumerate}

\(g\) ) \(d\) ) (i) 中的自同构的特征多项式为

\[
\mathrm{A} (\mathrm{T}) = (\mathrm{T} - 1) ^ {l} \prod_ {\alpha \in \mathrm{R}} (\mathrm{T} - \zeta^ {\alpha (H)}).
\]

自同构 \(s\) 在 \(e\) ) 中的特征多项式为

\[
\mathrm{B} (\mathrm{T}) = \left(\mathrm{T} ^ {h} - 1\right) ^ {l} \prod_ {i = 1} ^ {l} \left(\mathrm{T} - \zeta^ {m _ {i}}\right) \quad \left(m _ {i} \text { being   the   exponents   of } \mathrm{R}\right).
\]

（使用第VI章§1第11节的命题33 (iv)。）由关系 \(\mathrm{A(T)} = \mathrm{B(T)}\) 推出，对于所有 \(j \geq 1\)，满足该条件的 \(i\)（即 \(m_i \geq j\)）的个数等于满足条件 \(\alpha\in R_{+}\) 且 \(\alpha(H)=j\) 的元素个数；从而重新得到第VI章§4的习题6 c) 的结果。\(^{12}\)

\begin{enumerate}
\def\labelenumi{\arabic{enumi})}
\setcounter{enumi}{5}
\item
  假设 \(k = \mathbf{R}\) 或 \(\mathbf{C}\)。令 \(G\) 为李群 \(\mathrm{Aut}_0(\mathfrak{g})\)。证明元素 \(a \in G\) 是正则元（按第VII章§4第2节的意义）当且仅当它满足习题5 a) 的条件。
\item
  假设 g 是可分裂的。令 \(\mathrm{B}(\mathfrak{g})\) 为 g 的典则嘉当子代数的典则基（第3节，注2）。若 \(s \in \mathrm{Aut}(\mathfrak{g})\)，则 s 在 \(\mathrm{B}(\mathfrak{g})\) 上诱导一个置换；以 \(\operatorname{sgn}(s)\) 表示该置换的符号。证明 s 在 \(\bigwedge^{n} g\)（其中 \(n = \dim g\)）上的作用为
\end{enumerate}

\[
x \mapsto \operatorname{sgn} (s). x.
\]

¶8) 令 \(\bar{k}\) 为 k 的代数闭包，且 \(\bar{g} = \bar{k} \otimes_{k} g\)。Galois 群 \(\operatorname{Gal}(\bar{k}/k)\) 自然地作用于 \(\bar{g}\)、\(\bar{g}\) 的典则嘉当子代数（第3节，注2），以及其根系 \(\bar{R}\) 和典则基 \(\bar{B}\)。因而得到一个连续同态（即~具有开核）

\[
\pi : \operatorname{Gal} (\bar {k} / k) \to \operatorname{Aut} (\bar {\mathrm{R}}, \bar {\mathrm{B}}).
\]

证明 \(\mathfrak{g}\) 是可分裂的，当且仅当满足下列两个条件：

\begin{enumerate}
\def\labelenumi{(\roman{enumi})}
\item
  同态 \(\pi\) 是平凡的。
\item
  g 有一个 Borel 子代数 b。
\end{enumerate}

(证明，\(\mathfrak{h}\) 包含于 \(\mathfrak{b}\) 中的情形是分裂的，当且仅当 \(\pi\) 是平凡的。)

\begin{enumerate}
\def\labelenumi{\arabic{enumi})}
\setcounter{enumi}{8}
\tightlist
\item
  令 R 为既约根系，B 为 R 的一个基，且 \((\mathfrak{g}_{0}, \mathfrak{h}_{0}, \mathrm{B}, (X_{\alpha})_{\alpha \in \mathrm{B}})\) 为相应的带标架半单李代数（§4）。令 \(\bar{k}\) 为 k 的代数闭包，令 \(\rho : \operatorname{Gal}(\bar{k}/k) \to \operatorname{Aut}(\mathrm{R}, \mathrm{B})\) 为连续同态（参见~习题8）；若 \(\sigma \in \operatorname{Gal}(\bar{k}/k)\)，以 \(\rho_{\sigma}\) 表示 \(\bar{g}_{0} = \bar{k} \otimes_{k} g_{0}\) 的 k-线性自同构，使得 \(\rho_{\sigma}(X_{\alpha}) = X_{\rho(\sigma)\alpha}\)。另一方面，\(\operatorname{Gal}(\bar{k}/k)\) 在 \(\bar{k}\) 上的自然作用可以延拓为在 \(\bar{g}_{0}\) 上的作用。令 g 为 \(\bar{g}_{0}\) 的子集，由满足 \(\rho_{\sigma}(x) = \sigma^{-1}.x\) 对所有 \(\sigma \in \operatorname{Gal}(\bar{k}/k)\) 成立的元素 x 构成。
\end{enumerate}

\begin{enumerate}
\def\labelenumi{\alph{enumi})}
\item
  证明 \(\mathfrak{g}\) 是 \(k\)-李子代数，包含于 \(\bar{\mathfrak{g}}_0\)；并且 \(\mathfrak{g}\) 到 \(\bar{\mathfrak{g}}_0\) 的嵌入延拓为从 \(\bar{k} \otimes_k \mathfrak{g}\) 到 \(\bar{\mathfrak{g}}_0\) 的一个同构。特别地，\(\mathfrak{g}\) 是半单的。
\item
  令 \(\mathfrak{b}_0\) 为 \(\mathfrak{g}_0\) 的子代数，它由 \(\mathfrak{h}_0\) 和 \(X_{\alpha}\) 生成。置
\end{enumerate}

\[
\bar {\mathfrak {b}} _ {0} = \bar {k} \otimes_ {k} \mathfrak {b} _ {0}, \quad \mathfrak {h} = \mathfrak {g} \cap \bar {\mathfrak {b}} _ {0}, \quad \bar {\mathfrak {h}} _ {0} = \bar {k} \otimes_ {k} \mathfrak {h} _ {0}, \quad \mathfrak {h} = \mathfrak {g} \cap \bar {\mathfrak {h}} _ {0}.
\]

证明 \(\bar{k}\otimes_{k}b=\bar{b}_{0}\) 且 \(\bar{k}\otimes_{k}h=\bar{h}_{0}\)，从而 b 是 g 的 Borel 子代数，而 h 是包含于其中的嘉当子代数。

\begin{enumerate}
\def\labelenumi{\alph{enumi})}
\setcounter{enumi}{2}
\tightlist
\item
  证明与代数 g 相关联的同态 \(\pi\)（参见~习题8）等于 \(\rho\)。
\end{enumerate}

¶ 10) 令 h 为 g 的分裂嘉当子代数，且 \((X_{\alpha})_{\alpha\in\mathbb{R}}\) 为 \((\mathfrak{g},\mathfrak{h})\) 中的谢瓦莱系，参见~§2第4节。若 \(\alpha\in R\)，置

\[
\theta_ {\alpha} = e ^ {\mathrm{ad} X _ {\alpha}} e ^ {\mathrm{ad} X _ {- \alpha}} e ^ {\mathrm{ad} X _ {\alpha}},
\]

并以 \(\overline{W}\) 表示 \(\operatorname{Aut}_{e}(\mathfrak{g},\mathfrak{h})\) 中由 \(\theta_{\alpha}\) 生成的子群。

\begin{enumerate}
\def\labelenumi{\alph{enumi})}
\item
  证明 \(\varepsilon (\overline{\mathrm{W}}) = \mathrm{W}(\mathrm{R})\)。
\item
  令 \(s \in \overline{\mathrm{W}}\)，令 \(w = \varepsilon(s)\)。证明
\end{enumerate}

\[
s (X _ {\alpha}) = \pm X _ {w (\alpha)} \quad \text {   for   all   } \alpha \in \mathrm{R}.
\]

（使用§2的习题5。）

\begin{enumerate}
\def\labelenumi{\alph{enumi})}
\setcounter{enumi}{2}
\tightlist
\item
  令 M 为 \(\varepsilon : \overline{\mathrm{W}} \to \mathrm{W}(\mathrm{R})\) 的核。证明 M 包含于 \(f(\mathrm{T}_{\mathrm{Q}})\) 的子群，该子群由元素 \(f(\varphi)\) 组成，其中 \(\varphi^2 = 1\)。证明 M 包含由下式定义的元素 \(f(\varphi_{\alpha})\)，其中 \(\varphi_{\alpha}(\beta) = (-1)^{\langle \beta, \alpha^{\vee} \rangle}\)（注意 \(\theta_{\alpha}^{2} = f(\varphi_{\alpha})\)。）
\end{enumerate}

\(d)\) *令 \(\varphi\in\mathrm{Hom}(\mathbf{Q},\{\pm1\})\)。证明 \(f(\varphi)\) 属于 M，当且仅当 \(\varphi\) 延拓为一个从 P 到 \(\{\pm1\}\) 的同态。（充分性由 M 包含 \(f(\varphi_{\alpha})\) 这一事实得到。为证明必要性，化归到 \(k=\mathbf{Q}\) 的情形，并使用 M 包含于 \(f(\mathrm{T}_{\mathrm{Q}})\cap\mathrm{Aut}_{e}(\mathfrak{g})=\mathrm{Im}(\mathrm{T}_{\mathrm{P}})\) 的事实，参见~§7，习题26 \(d)\)。推出 M 同构于群 \(\mathrm{Q}/(\mathrm{Q}\cap2\mathrm{P}).^{13}\) 的对偶。

\begin{enumerate}
\def\labelenumi{\arabic{enumi})}
\setcounter{enumi}{10}
\tightlist
\item
  按第2节的记号，假设 \(k\) 是非离散且局部紧的，因而同构于 \(\mathbf{R}\)、\(\mathbf{C}\) 或 \(\mathbf{Q}_p\) 的有限扩张（《交换代数》，第VI章§9第3节）。对所有 \(n \geq 1\)，商 \(k^* / k^{*n}\) 是有限的（超度量情形参见~《交换代数》第VI章§9的习题3）。推出商 \(T_Q / \text{Im}(T_P)\) 和 \(\text{Aut}(\mathfrak{g}) / \text{Aut}_e(\mathfrak{g})\) 是有限的。
\end{enumerate}

当 \(k = \mathbf{R}\) 时，证明 \(\mathrm{T_Q / Im(T_P)}\) 同构于 \(\mathbf{F}_{2}\)-向量空间 \((\mathrm{Q} \cap 2\mathrm{P}) / 2\mathrm{Q}\) 的对偶。当 \(k = \mathbf{C}\) 时，有 \(\mathrm{T_Q = Im(T_P)}\)；这是李群 \(\operatorname{Aut}(\mathfrak{g})\) 中李代数为 \(\mathfrak{h}\) 的整子群。

¶ 12) 令 \(\mathfrak{h}\) 为 \(\mathfrak{g}\) 的分裂嘉当子代数，A 为 \(\mathfrak{h}\) 的子集，且 \(s \in \operatorname{Aut}(\mathfrak{g})\) 满足 \(s\mathrm{A} = \mathrm{A}\)。证明存在 \(t \in \operatorname{Aut}(\mathfrak{g}, \mathfrak{h})\)，使得 \(t|\mathrm{A} = s|\mathrm{A}\) 且 \(ts^{-1} \in \operatorname{Aut}_e(\mathfrak{g})\)。（令 \(\mathfrak{a}\) 为 A 在 \(\mathfrak{g}\) 中的交换子；这是 \(\mathfrak{g}\) 的约化子代数，其中 \(sh\) 和 \(\mathfrak{h}\) 是分裂嘉当子代数；推出存在 \(u \in \operatorname{Aut}_e(\mathfrak{a})\)，使得 \(ush = \mathfrak{h}\)；证明存在 \(v \in \operatorname{Aut}_e(\mathfrak{g})\) 延拓 \(u\)，使得 \(v|\mathrm{A} = \mathrm{Id}_{\mathrm{A}}\)；取 \(t = vs.\)）推出，如果 \(s \in \operatorname{Aut}_0(\mathfrak{g})\)，则存在 \(w \in \mathrm{W}(\mathrm{R})\)，使得 \(w|\mathrm{A} = s|\mathrm{A}\)。

¶13) 令 \((\mathfrak{g},\mathfrak{h},\mathrm{B},(X_{\alpha})_{\alpha\in\mathrm{B}})\) 为带标架半单李代数，R 为 \((\mathfrak{g},\mathfrak{h})\) 的根系，\(\Delta\) 为相应的邓金图，且 \(\varPhi\) 为下式的一个子群

\[
\operatorname{Aut} (\mathrm{R}, \mathrm{B}) = \operatorname{Aut} (\Delta).
\]

若 \(s \in \varPhi\)，则按以下条件将 \(s\) 延拓为 \(\mathfrak{g}\) 的一个自同构

\[
s (X _ {\alpha}) = X _ {s \alpha}, \quad \text { and } \quad s (H _ {\alpha}) = H _ {s \alpha} \quad \text { for   all } \alpha \in \mathrm{B}, \text { cf.   Prop. } 1.
\]

于是将 \(\varPhi\) 与 \(\operatorname{Aut}(\mathfrak{g},\mathfrak{h})\) 的一个子群相识别；以 \(\tilde{\mathfrak{g}}\)（分别地，\(\tilde{\mathfrak{h}}\)）表示 g（分别地，h）中在 \(\varPhi\) 下不变的元素所成的子代数。

\begin{enumerate}
\def\labelenumi{\alph{enumi})}
\tightlist
\item
  令 \(\alpha \in \mathbf{B}\)，令 \(\mathbf{X} = \varPhi.\alpha\)。利用第VI章中的图版，证明只有以下两种情形可能：
\end{enumerate}

\begin{enumerate}
\def\labelenumi{(\roman{enumi})}
\item
  \(X\) 中异于 \(\alpha\) 的每个元素都与 \(\alpha\) 正交；
\item
  存在唯一元素 \(\alpha'\)，它属于 \(X - \{\alpha\}\)，且不与 \(\alpha\) 正交，并且 \(n(\alpha, \alpha') = n(\alpha', \alpha) = -1\)。
\end{enumerate}

\begin{enumerate}
\def\labelenumi{\alph{enumi})}
\setcounter{enumi}{1}
\tightlist
\item
  令 i 为限制映射 \(h^{*} \to \tilde{h}^{*}\)，令 \(\tilde{\mathbf{B}} = i(\mathbf{B})\)；映射 \(\mathbf{B} \to \tilde{\mathbf{B}}\) 将 \(\tilde{B}\) 与 \(B/\varPhi\) 相识别。证明 \(\tilde{g}\) 是半单李代数，\(\tilde{h}\) 是其中的分裂嘉当子代数，且 \(\tilde{B}\) 是 \(\mathrm{R}(\tilde{\mathfrak{g}}, \tilde{\mathfrak{h}})\) 的一个基。（注意 \(\tilde{B}\) 包含于 \(\mathrm{R}(\tilde{\mathfrak{g}}, \tilde{\mathfrak{h}})\) 中，且 \(\mathrm{R}(\tilde{\mathfrak{g}}, \tilde{\mathfrak{h}})\) 的每个元素都是 \(\tilde{B}\) 中元素的具有相同符号的整数系数线性组合。）若 \(\tilde{\alpha} \in \tilde{B}\)，相应的逆根 \(H_{\tilde{\alpha}} \in \tilde{h}\) 由下式给出
\end{enumerate}

\[
\begin{array}{l} H _ {\tilde {\alpha}} = \sum_ {i (\alpha) = \tilde {\alpha}} H _ {\alpha} \quad \text {in case (i) of a)} \\ H _ {\tilde {\alpha}} = 2 \sum_ {i (\alpha) = \tilde {\alpha}} H _ {\alpha} \quad \text {in case (ii) of a)}, \end{array}
\]

其中求和遍历 \(\alpha \in \mathrm{B}\) 中满足 \(i(\alpha) = \tilde{\alpha}\) 的元素。若 \(\beta \in \mathrm{B}\) 的像为 \(\tilde{\beta} \in \tilde{\mathrm{B}}\)，则

\[
  n (\tilde {\beta}, \tilde {\alpha}) = \sum_ {i (\alpha) = \tilde {\alpha}} n (\beta , \alpha) \quad \text { in   case   (i) }
\]

\[
n (\tilde {\beta}, \tilde {\alpha}) = 2 \sum_ {i (\alpha) = \tilde {\alpha}} n (\beta , \alpha) \quad \text { in   case   (ii) }.
\]

推出 \(\mathrm{R}(\tilde{\mathfrak{g}},\tilde{\mathfrak{h}})\) 的邓金图如何由二元组 \((\varDelta,\varPhi)\) 确定。

\begin{enumerate}
\def\labelenumi{\alph{enumi})}
\setcounter{enumi}{2}
\tightlist
\item
  证明若 \(\mathfrak{g}\) 是单的，则 \(\tilde{\mathfrak{g}}\) 也是单的。
\end{enumerate}

若 \(\mathfrak{g}\) 属于 \(A_l\) 型，\(l \geq 2\)，且 \(\varPhi\) 的阶为2，则 \(\tilde{\mathfrak{g}}\) 属于 \(B_{l/2}\) 型（当 \(l\) 为偶数时），属于 \(C_{(l+1)/2}\) 型（当 \(l\) 为奇数时）。

若 \(\mathfrak{g}\) 属于 \(D_l\) 型，\(l \geq 4\)，且 \(\varPhi\) 的阶为2，则 \(\tilde{\mathfrak{g}}\) 属于 \(B_{l-1}\) 型。

若 \(\mathfrak{g}\) 属于 \(\mathrm{D}_4\) 型，且 \(\varPhi\) 的阶为3或6，则 \(\tilde{\mathfrak{g}}\) 属于 \(\mathrm{G}_2\) 型。

若 g 属于 \(E_{6}\) 型，且 \(\Phi\) 的阶为2，则 \(\tilde{g}\) 属于 \(F_{4}\) 型。

